\section{总结与延伸}
本讲围绕 inference-time compute 展开：先让提示释放 token，再用 sampling/search 扩展 solution space，最后靠 trace + self-correction 修正结果。三段可以按需组合，在高质量 reasoning 任务中显著提升 accuracy 与 reliability。

我们还梳理了 compute scheduling（三角调度）、self-consistency/search 的协同机制、stepwise verifier 与 Tree-of-Thought 的控制 loop，以及 benchmark 的数据解读，形成一张从 prompt 到 deployment 的完整路线图。

\begin{importantbox}{确认落地的要点}
1. 在工程中持续追踪 sampling、search、trace 三类指标；2. 把 benchmark 曲线（例如 Arc AGI）作为 budget gating 的参考；3. 把 deployment SLO、alert、budget instrumentation 写入 CI/CD pipeline，确保 inference-time compute 不会超出可控范围。
\end{importantbox}

\begin{table}[H]
\centering
\caption{推理时计算的核心套路}
\begin{tabular}{p{0.28\textwidth}p{0.32\textwidth}p{0.32\textwidth}}
\toprule
策略 & 关键收益 & 典型代价 \\
\midrule
扩展 CoT 预算 & 让模型在难题上逼近 human-level accuracy & 更多 token、更多 latency、prompt 设计更敏感 \\
多候选采样与 self-consistency & 从 mistakes 中恢复，适用于没有 verifier 的场景 & 每条 candidate 都要 compute，需 voter pool 与一致性表征 \\
Tree-of-Thought 与 step verifier & 早期淘汰 bad paths、与 verifier 配合提升可解释性 & branching 数量激增，需要设计 quality function \\
Trace feedback + self-correction & Line-by-line 评分可逐步提升 accuracy & 无 oracle 会导致 feedback drift、需要 prompt guard rails \\
\bottomrule
\end{tabular}
\end{table}

\subsection{本章小结}
推理时技术需把提示、search、trace 合成 pipeline，通过预算调度与指标回路让每个模块在不同阶段发挥最大效益。

\subsection{进一步阅读}
\begin{itemize}
\item Wang 等人《Self-Consistency improves chain of thought reasoning》阐释了 sampling + majority vote 的 scaling 效果。
\item Zhou 等人《Tree-of-Thought》提出通过树状探索把 search 与 step evaluation 结合。
\item 《Trace-Enhanced LLM Safety》等 trace feedback 论文说明 natural language trace 如何辅助 verifier。
\item DeepMind 的 AlphaProof/AlphaGeometry/AlphaCode 系列展示了 inference-time RL + search + trace 的组合在复杂 reasoning 任务中的成果。
\end{itemize}

\end{document}
